% RESEARCH_DRAFT_SHA256: a57dbb879e3d2341fa853279d2b3d41df69c66e7c48f8a6518c264153c1a9941
% Layout and source-reference wrappers expanded for the web renderer.
\documentclass[11pt]{article}
\usepackage[margin=25mm]{geometry}
\usepackage{amsthm}
\usepackage{fontspec,unicode-math}
\setmainfont{Latin Modern Roman}
\setmathfont{Latin Modern Math}
\usepackage{amsmath,amssymb,mathtools,xurl,hyperref}
\hypersetup{colorlinks=true,urlcolor=blue,linkcolor=blue}
\setcounter{secnumdepth}{0}
\setlength{\parindent}{0pt}
\setlength{\parskip}{5pt}
\setlength{\emergencystretch}{3em}
\newcommand{\sref}[2]{\hyperlink{source:#1:#2}{[S#2]}}
\newcommand{\cataloguescope}[1]{\paragraph{Source record.}#1}
\newcommand{\floor}[1]{\left\lfloor#1\right\rfloor}
\newtheorem{proposition}{Proposition}
\newtheorem{lemma}[proposition]{Lemma}
\begin{document}
\paragraph{Research manuscript.} OpenAI \texttt{gpt-6-astra}, reasoning \texttt{ultra}; 11 October 2026.
\paragraph{Preserved source body.} The complete catalogue mathematical text and sources follow with layout and source links expanded, including source qualifications and its historical preparation notes. New research begins at the section headed Research attempt.

% UnsolvedMath ID 20002185; source catalogue code AIM-INFRASTRUCTURE-0085
\section{A degree-slice fast path for basis}\label{20002185}
{\small UnsolvedMath ID 20002185 · Computer Science}
\subsection{Definitions and mathematical statement}

\paragraph{Shared notation.}Standard number systems are $\mathbb N,\mathbb Z,\mathbb Q,\mathbb R,\mathbb C$. Specialized notation and index ranges follow the cited source and the definitions below.


For a graded module $M=\bigoplus_d M_d$ over a polynomial ring with coefficient field $k$, a basis of a finite-dimensional component $M_d$ is a linearly independent spanning set over $k$. Macaulay2 basis can request a particular degree, a degree range, or a generating set over a coefficient ring. Such a result is different from asserting that $M$ is free over its polynomial ring.
\paragraph{Record scope and qualifications.}
The catalogue title is retained as metadata. Named certificates, contracts, bounds or frameworks in that title are not asserted as results of the original source extract.
The referenced issue and a version-specific slow example are not included. Faster and without breaking anything require an input family, correctness contract, regression scope, and comparison implementation; no bound is invented.

Catalogue screening limitation: Software performance bug linked to a GitHub issue; the source supplies no mathematical input family or complexity objective.

The source snapshot labels this record partially solved; this is catalogue metadata rather than independent verification of a mathematical result or the current completion of a project.
\paragraph{Statement recorded in the supplied source.}
Other internal improvements

The current "basis" method can be quite slow in situations where it should be nearly instantaneous (see Github issue \#2899).

How can we implement a faster basis calculation without breaking anything? \href{http://aimpl.org/macaulay2efie/1/}{[S2]}
\paragraph{Source status.}
The supplied catalogue labels this source record partially solved; that metadata was checked against the supplied snapshot on 2026-10-11. The label is preserved for provenance and is not an independent verification of present open status or mathematical correctness.
\subsection{Short English statement}
Improve the performance of Macaulay2 basis calculations while preserving their documented behavior.
\subsection{Sources}
\begin{itemize}
\item[S1] Ulam AI and the UnsolvedMath Contributors, UnsolvedMath: A Curated Collection of Open Mathematics Problems, record 20002185 (AIM-INFRASTRUCTURE-0085). Catalogue curation and metadata: CC BY 4.0. The source extract is retained; definitions, English summaries and qualifications are editorial additions.
\url{https://huggingface.co/datasets/ulamai/UnsolvedMath}.
\item[S2] AIM Workshop Problem Lists, source linked by catalogue record AIM-INFRASTRUCTURE-0085.
\url{http://aimpl.org/macaulay2efie/1/}.
\item[S3] Official Macaulay2 documentation, basis: basis or generating set of all or part of a ring, ideal or module.
\url{https://macaulay2.com/doc/Macaulay2/share/doc/Macaulay2/Macaulay2Doc/html/_basis.html}.
\end{itemize}
\paragraph{Preparation and review.}
The original source and UnsolvedMath contributors retain attribution for the source extract. Codex (GPT-6; exact serving revision not exposed) assisted with definitions, English summaries, scope notes and formatting. No new human mathematical authorship or independently verified solution is asserted. Permanent source ID 20002185 and code AIM-INFRASTRUCTURE-0085 checked against the complete supplied registry and current repository. Distinct upstream records, including repeated questions and broad or damaged extracts, are retained for complete coverage of the source CS category. A repeated record is not a new independent conjecture. Complete mathematical-equivalence and current-literature review has not been performed. Preparation and final review follow contribution prompts C and D, with source gaps disclosed. Structural and rendering checks do not establish a new theorem or open status.
\label{end:20002185}

\section{Research attempt}
\subsection{Literature, source context and scope}
The complete catalogue body above retains the original agenda, source status, and historical preparation. The AIMPL problem page could not be retrieved. The supplied complete source transcription was read, and the accessible primary AIM workshop announcement was checked for context. The announcement is not a substitute for verification of the inaccessible problem page. The catalogue's earlier proposed frameworks and research summaries are context, not established solutions credited to this attempt. The following is an explicitly limited investigation of one route; it neither substitutes a formalized toy model for the entire source agenda nor claims empirical or deployed-software results. Elementary auxiliary facts are not new original-target progress.

\subsection{A degree-slice prototype with a proved basis contract}
Official documentation distinguishes vector-space bases over a coefficient field from generating sets over more general coefficient rings. We restrict to a field $K$ and $R=K[x,y]/(x^a,y^b)$, where $a,b\ge1$, with the usual total grading. The goal of this prototype is the homogeneous degree-$D$ vector-space basis for $D\ge0$, not the general overloaded Macaulay2 basis method.
\begin{proposition}
Set $\ell=\max(0,D-b+1)$ and $u=\min(D,a-1)$. The monomials
\[
\{x^i y^{D-i}:\ell\le i\le u\}
\]
form a basis of $R_D$, of cardinality $\max(0,u-\ell+1)$.
\end{proposition}
\begin{proof}
The ideal $(x^a,y^b)$ is spanned over $K$ by monomials with $x$ exponent at least $a$ or $y$ exponent at least $b$: multiples of either generator have this property, and every such monomial is a multiple of a generator. Since distinct monomials form a basis of the polynomial ring, the complementary monomials with $0\le i<a$ and $0\le j<b$ give a basis of the quotient: they span after deleting ideal monomials, and any linear combination in the ideal must have all complementary coefficients zero. Imposing $i+j=D$ yields $j=D-i\ge0$, $i\ge0$, $i\le a-1$, and $D-i\le b-1$, exactly $\ell\le i\le u$. Counting integers proves the formula.
\end{proof}
A direct generator computes the two bounds and emits only that interval of exponent pairs. For $(a,b,D)=(4,3,4)$ it returns $(i,D-i)=(2,2),(3,1)$, namely $x^2y^2,x^3y$. It never constructs the whole quotient basis. The formula requires only a constant number of integer bound operations plus one output operation per returned pair; this is not a bit-complexity or Macaulay2 runtime claim.
\subsection{Verification and remaining gaps}
The generator was compared with exhaustive monomial enumeration for $1\le a,b\le8$ and $0\le D\le16$, all 1088 cases agreeing. The referenced issue 2899 could not be retrieved, so its exact failing case and current fix status were not verified. No M2 patch, dispatch compatibility, cache behavior, general module basis, or coefficient-ring case was addressed. The prototype supplies only a familiar special-case fast path.

\section{References and source attribution}
Primary source: \url{http://aimpl.org/macaulay2efie/1/}. Attempted primary AIMPL page retrieval; failed with timeout or gateway error. The supplied full source transcription was read; the inaccessible live page was not independently inspected.

\url{https://aimath.org/pastworkshops/macaulay2efie.html}. Read the primary September 2023 workshop announcement.

\url{https://aimath.org/pastworkshops/macaulay2efierep.pdf}. Read the complete three-page primary workshop report. Existing threads/tasks, MPI work and parallelApply are acknowledged; no absence of all existing parallelism is inferred.

\url{https://macaulay2.com/doc/Macaulay2/share/doc/Macaulay2/Macaulay2Doc/html/_basis.html}. Read official basis documentation, including graded components and the distinction between coefficient-field bases and more general generating sets.

\url{https://github.com/Macaulay2/M2/issues/2899}. Attempted primary issue retrieval and API fallback; both failed. Did not inspect its actual discussion or claim a current fix.

All retrievals were checked on 11 October 2026. Original contributor and catalogue attribution remains in the preserved Sources section. The elementary arguments below are self-contained; their exposition is not a claim of first discovery.


\paragraph{Integration review.} Codex (GPT-6; exact serving revision not exposed) reviewed the written arguments and expanded layout and source-reference wrappers. Repository metadata and the production web renderer were checked using Python 3.13 and JavaScript with MathJax 3.2.2. This model-assisted review is not independent human proof verification; no PDF compilation or formal proof checking is asserted.

\section{Completion Estimate}
\noindent\textbf{COMPLETION: 0\%}
No substantiated new progress toward the entire original target is claimed. The worked results and failed route are delimited above; a known result, illustrative model, or auxiliary obstruction is not counted as original-target completion.

\section{Final status and attribution}
\textbf{Attempt status: unresolved.} This describes this research attempt, not a new ruling on the source catalogue's status or on external literature claims.
The mathematics and this self-audit were generated by OpenAI \texttt{gpt-6-astra}, reasoning \texttt{ultra}, on 11 October 2026. Codex (GPT-6; exact serving revision not exposed) orchestrated the batch and contributed formatting/repository review. No human mathematical authorship, independent human proof checking, or formal verification is claimed. The historical preparation paragraphs in the reproduced catalogue body concern that body only.

Contribution prompt A was applied to the complete definition and supplied source record, followed by the adversarial checks of prompt B. Local reproduction files are \texttt{batch013/checks.py}, \texttt{computations.json}, and the per-ID audit JSON. They were not public links at preparation time. Numerical checks are finite and do not prove unrestricted results. No LaTeX engine was found locally; compilation was not performed by this research agent.

\end{document}
